% GENERATED FILE. Edit canonical lessons, then run npm run generate:books.
% canonical-source-hash: fnv1a64:0493a6d83829f47a
% canonical-lessons: TE-C336-pukaru, TE-C336-vadanti, TE-C336-atisayokti, TE-C336-kalakalam, TE-C336-velladincu

\chapter{A Rumour Spreads}
\label{ch:te-a-rumour-spreads}
\hlchaptermodality{336}

\begin{chapteropening}
\textbf{By the end of this chapter:} \emph{I can say a rumour, hearsay, an exaggeration, an uproar and to disclose.}

Complete the last lesson of chapter 336: I can say a rumour, hearsay, an exaggeration, an uproar and to disclose.
\end{chapteropening}

\section[\texorpdfstring{\te{పుకారు} (pukāru)}{పుకారు (pukāru)}]{\te{పుకారు} (pukāru) — a rumour}

\label{lesson:TE-C336-pukaru}



\begin{quote}
Before the new one: say the Telugu for to enthrone, then the Telugu for a reign.
\end{quote}

\subsection*{You'll want to know: \te{పుకారు}}
\textbf{\te{పుకారు}} — \emph{pukāru} — \textquotedblleft{}a rumour\textquotedblright{}.

\begin{quote}
\te{బజారులో} \te{ఒక} \te{పుకారు} \te{వినిపించింది}.
\end{quote}

\begin{quote}
\emph{bajārulō oka pukāru vinipiñcindi.}
\end{quote}

\begin{quote}
A rumour was going round the market.
\end{quote}

\te{పుకారు} \te{పుట్టించు}, to start a rumour.

\begin{cousinweb}
\textbf{\te{పుకారు}} (\emph{pukāru}) is the everyday Telugu word for a rumour; its origin is not settled.
\end{cousinweb}

\begin{grammarlens}[title={What you've built}]
One more word for telling what happened.
\end{grammarlens}

\subsection*{Your turn}
\begin{itemize}
\raggedright
  \item \emph{Say it:} \emph{pukāru}
  \item \emph{Say it:} \emph{pukāru}, once more
  \item \emph{Say it:} \emph{bajārulō oka pukāru vinipiñcindi.} — A rumour was going round the market.
  \item \emph{Recall:} say the Telugu for to enthrone, then the Telugu for a reign, then say \emph{pukāru} again
\end{itemize}

\begin{tcolorbox}[breakable,colback=teal!4,colframe=teal!35!black,title={Before you move on}]
What does \te{పుకారు} mean? (\textquotedblleft{}A rumour\textquotedblright{}.) Say it once more.
\end{tcolorbox}
\section[\texorpdfstring{\te{వదంతి} (vadanti)}{వదంతి (vadanti)}]{\te{వదంతి} (vadanti) — hearsay; a rumour}

\label{lesson:TE-C336-vadanti}



\begin{quote}
Before the new one: say the Telugu for a reign, then the Telugu for a rumour.
\end{quote}

\subsection*{You'll want to know: \te{వదంతి}}
\textbf{\te{వదంతి}} — \emph{vadanti} — \textquotedblleft{}hearsay; a rumour\textquotedblright{}.

\begin{quote}
\te{అది} \te{నిజం} \te{కాదు}, \te{వదంతి} \te{మాత్రమే}.
\end{quote}

\begin{quote}
\emph{adi nijaṁ kādu, vadanti mātramē.}
\end{quote}

\begin{quote}
It is not true, only hearsay.
\end{quote}

Close to \te{పుకారు}; news writers prefer \te{వదంతి}.

\begin{cousinweb}
\textbf{\te{వదంతి}} (\emph{vadanti}) goes back to Sanskrit \emph{kiṃvadantī}, \textquotedblleft{}a common saying, a rumour\textquotedblright{}, literally \textquotedblleft{}what do they say\textquotedblright{} (Monier-Williams), which Monier-Williams also lists in the short form \emph{vadantī}.
\end{cousinweb}

\begin{grammarlens}[title={What you've built}]
One more word for telling what happened.
\end{grammarlens}

\subsection*{Your turn}
\begin{itemize}
\raggedright
  \item \emph{Say it:} \emph{vadanti}
  \item \emph{Say it:} \emph{vadanti}, once more
  \item \emph{Say it:} \emph{adi nijaṁ kādu, vadanti mātramē.} — It is not true, only hearsay.
  \item \emph{Recall:} say the Telugu for a reign, then the Telugu for a rumour, then say \emph{vadanti} again
\end{itemize}

\begin{tcolorbox}[breakable,colback=teal!4,colframe=teal!35!black,title={Before you move on}]
What does \te{వదంతి} mean? (\textquotedblleft{}Hearsay; a rumour\textquotedblright{}.) Say it once more.
\end{tcolorbox}
\section[\texorpdfstring{\te{అతిశయోక్తి} (atiśayōkti)}{అతిశయోక్తి (atiśayōkti)}]{\te{అతిశయోక్తి} (atiśayōkti) — an exaggeration}

\label{lesson:TE-C336-atisayokti}



\begin{quote}
Before the new one: say the Telugu for a rumour, then the Telugu for hearsay.
\end{quote}

\subsection*{You'll want to know: \te{అతిశయోక్తి}}
\textbf{\te{అతిశయోక్తి}} — \emph{atiśayōkti} — \textquotedblleft{}an exaggeration\textquotedblright{}.

\begin{quote}
\te{అతని} \te{మాటల్లో} \te{చాలా} \te{అతిశయోక్తి} \te{ఉంది}.
\end{quote}

\begin{quote}
\emph{atani māṭallō cālā atiśayōkti undi.}
\end{quote}

\begin{quote}
There is a lot of exaggeration in his words.
\end{quote}

\te{అతిశయ}, excess, and \te{ఉక్తి}, a saying.

\begin{cousinweb}
\textbf{\te{అతిశయోక్తి}} (\emph{atiśayōkti}) is borrowed from Sanskrit \emph{atiśayokti}, \textquotedblleft{}hyperbolical language\textquotedblright{} (Monier-Williams), from \emph{atiśaya}, \textquotedblleft{}excess\textquotedblright{}, and \emph{ukti}, \textquotedblleft{}a saying\textquotedblright{}.
\end{cousinweb}

\begin{grammarlens}[title={What you've built}]
One more word for telling what happened.
\end{grammarlens}

\subsection*{Your turn}
\begin{itemize}
\raggedright
  \item \emph{Say it:} \emph{atiśayōkti}
  \item \emph{Say it:} \emph{atiśayōkti}, once more
  \item \emph{Say it:} \emph{atani māṭallō cālā atiśayōkti undi.} — There is a lot of exaggeration in his words.
  \item \emph{Recall:} say the Telugu for a rumour, then the Telugu for hearsay, then say \emph{atiśayōkti} again
\end{itemize}

\begin{tcolorbox}[breakable,colback=teal!4,colframe=teal!35!black,title={Before you move on}]
What does \te{అతిశయోక్తి} mean? (\textquotedblleft{}An exaggeration\textquotedblright{}.) Say it once more.
\end{tcolorbox}
\section[\texorpdfstring{\te{కలకలం} (kalakalaṁ)}{కలకలం (kalakalaṁ)}]{\te{కలకలం} (kalakalaṁ) — an uproar; a stir}

\label{lesson:TE-C336-kalakalam}



\begin{quote}
Before the new one: say the Telugu for hearsay, then the Telugu for an exaggeration.
\end{quote}

\subsection*{You'll want to know: \te{కలకలం}}
\textbf{\te{కలకలం}} — \emph{kalakalaṁ} — \textquotedblleft{}an uproar; a stir\textquotedblright{}.

\begin{quote}
\te{ఆ} \te{వార్త} \te{పట్టణంలో} \te{కలకలం} \te{రేపింది}.
\end{quote}

\begin{quote}
\emph{ā vārta paṭṭaṇaṁlō kalakalaṁ rēpindi.}
\end{quote}

\begin{quote}
The news caused an uproar in the town.
\end{quote}

\te{కలకలం} \te{రేపు}, to cause a stir, is the headline phrase.

\begin{cousinweb}
\textbf{\te{కలకలం}} (\emph{kalakalaṁ}) is borrowed from Sanskrit \emph{kalakala}, \textquotedblleft{}any confused noise, the murmur of a crowd\textquotedblright{} (Monier-Williams).
\end{cousinweb}

\begin{grammarlens}[title={What you've built}]
One more word for telling what happened.
\end{grammarlens}

\subsection*{Your turn}
\begin{itemize}
\raggedright
  \item \emph{Say it:} \emph{kalakalaṁ}
  \item \emph{Say it:} \emph{kalakalaṁ}, once more
  \item \emph{Say it:} \emph{ā vārta paṭṭaṇaṁlō kalakalaṁ rēpindi.} — The news caused an uproar in the town.
  \item \emph{Recall:} say the Telugu for hearsay, then the Telugu for an exaggeration, then say \emph{kalakalaṁ} again
\end{itemize}

\begin{tcolorbox}[breakable,colback=teal!4,colframe=teal!35!black,title={Before you move on}]
What does \te{కలకలం} mean? (\textquotedblleft{}An uproar; a stir\textquotedblright{}.) Say it once more.
\end{tcolorbox}
\section[\texorpdfstring{\te{వెల్లడించు} (vellaḍiñcu)}{వెల్లడించు (vellaḍiñcu)}]{\te{వెల్లడించు} (vellaḍiñcu) — to disclose; to reveal}

\label{lesson:TE-C336-velladincu}



\begin{quote}
Before the new one: say the Telugu for an exaggeration, then the Telugu for an uproar.
\end{quote}

\subsection*{You'll want to know: \te{వెల్లడించు}}
\textbf{\te{వెల్లడించు}} — \emph{vellaḍiñcu} — \textquotedblleft{}to disclose; to reveal\textquotedblright{}.

\begin{quote}
\te{పోలీసులు} \te{దొంగ} \te{పేరు} \te{వెల్లడించారు}.
\end{quote}

\begin{quote}
\emph{pōlīsulu doṅga pēru vellaḍiñcāru.}
\end{quote}

\begin{quote}
The police disclosed the thief's name.
\end{quote}

Past (I …): \te{వెల్లడించాను}, \emph{vellaḍiñcānu}.

What is disclosed takes no ending; to whom takes -\te{కి}.

\begin{cousinweb}
\textbf{\te{వెల్లడించు}} (\emph{vellaḍiñcu}) is native Dravidian: DEDR 5498 lists Telugu \emph{vellaḍi}, \textquotedblleft{}open place, publicity, openness\textquotedblright{}, under the word for \textquotedblleft{}outside\textquotedblright{}, with Tamil \emph{veḷi}; to disclose is to bring into the open.
\end{cousinweb}

\begin{grammarlens}[title={What you've built}]
One more word for telling what happened.
\end{grammarlens}

\subsection*{Your turn}
\begin{itemize}
\raggedright
  \item \emph{Say it:} \emph{vellaḍiñcu}
  \item \emph{Say it:} \emph{vellaḍiñcu}, once more
  \item \emph{Say it:} \emph{pōlīsulu doṅga pēru vellaḍiñcāru.} — The police disclosed the thief's name.
  \item \emph{Recall:} say the Telugu for an exaggeration, then the Telugu for an uproar, then say \emph{vellaḍiñcu} again
\end{itemize}

\begin{tcolorbox}[breakable,colback=teal!4,colframe=teal!35!black,title={Before you move on}]
What does \te{వెల్లడించు} mean? (\textquotedblleft{}To disclose; to reveal\textquotedblright{}.) Say it once more.
\end{tcolorbox}
